\section*{Limitations}

\paragraph{Quantization and serving.} Q checkpoints are fake-quantized: dequantized W4 weights with BF16 compute on a batch-invariant serving path. Deployed low-bit kernels change the numerics and would add serving variation that we deliberately remove, so our results describe the effect of the weight change, not of a production kernel. We study two Qwen3 models and two W4 methods, one primary checkpoint each; checkpoint-to-checkpoint variation within a method is not covered by the fork results.

\paragraph{One simulator, one benchmark.} All episodes use one \tautwo{} user simulator (Qwen3-30B-A3B, greedy) and three \tautwo{} domains. Continuation effects include the simulator's responses, so a different simulator may change $C$ and $I$. Retail and airline are scored by DB success, which is not \tautwo{}'s official reward; official rewards are reported separately, and telecom uses its official deterministic task-basis success.

\paragraph{Blocks from a researcher-defined perturbation.} Replicate blocks come from prompt nonces, an input perturbation we chose; they are not draws of a natural deployment distribution. Conclusions for the fixed task pool and for new tasks are reported separately. In a development run of the paired success analysis with stand-in pairs (not a result), the task-clustered interval for new tasks had a half-width of about $\pm$4\pp{} for the all-model pool, so claims at a 3\pp{} resolution about new tasks are unlikely with the current task counts; more independent tasks, not more nonces, would narrow it.

\paragraph{Scope of the fork.} Only the first semantic action divergence per task and block is intervened on; later divergences, and differences in free text only (which can carry commitments or authorisation requests), are outside the primary scope. The causal reading of the branches depends on the replay acceptance gate; if restoration cannot be verified, the paper reports measurements only. Mechanism labels depend on the tool-type table, on rule-based checks and on an LLM labeller; about 80 human-reviewed states validate them on an enriched sample, which does not give whole-data accuracy.

\paragraph{Small confirmation set.} The telecom holdout has 34 tasks. Few events there limit the evidence for any mechanism and do not show its absence; non-inferiority at 3\pp{} is not equivalence.

\paragraph{Retired lines, reported transparently.} This project previously pursued a ``checkpoint lottery'' (same method, different calibration draw, different success) and an account in which W4 agents matched BF16 success while behaving differently. The first was tested in a pre-registered confirmation and the exploratory results did not replicate; this is not evidence that checkpoint effects are absent, nor of BF16/W4 success equivalence (Appendix~\ref{app:prereg}). For the second, success equivalence was never established, and the pre-registered fingerprint test turned out to have an invalid permutation null; after correction, no fingerprint was detected for W4 quantizer checkpoints (Appendices~\ref{app:audit}, \ref{app:fingerprint}). Both are kept in the appendix as background, and neither supports a claim in the main text.
